Agricultural subsurface (tile) drainage is one of the foundational
technologies of modern crop production: by lowering the water table on
poorly drained land it secures trafficability, root aeration, and yield
across vast areas of the temperate and irrigated world. Yet the same
engineered conduits that make these soils farmable also short-circuit the
landscape's natural buffering, delivering nitrate, phosphorus, salts, and
an expanding list of other constituents rapidly to ditches, streams, and
ultimately to coastal waters. Decades of field experiments, monitoring
campaigns, and modelling have produced a large and detailed literature on
how drained systems behave and on the practices---nutrient management,
controlled drainage, cover crops, edge-of-field treatment, salinity
reclamation---that might reconcile productivity with water quality. The
distinctive feature of that literature, however, is that its most secure
and most frequently repeated statements are not about what is known but
about what remains unknown. This review is organised around that
observation.

Read across the field, the recurring signal is a precondition rather than
a conclusion. Headline results are drawn disproportionately from short
records, single sites, and narrow agro-ecological windows, and their
authors caution explicitly against premature generalisation. The blunt
appeals are pervasive---that ``more research is needed''
\parencite{SHALIT1995}, that ``the complexity of this subject requires
more research'' \parencite{Malone2014}, that ``further research in this
area is needed'' \parencite{Woo2023}, and that ``further evaluations are
warranted before any concrete conclusions could be drawn''
\parencite{GUPTA1993}---and they mark genuine gaps in evidence rather than
rhetorical closings. \textcite{Rogger2017} pitch the same concern at the
level of the whole discipline, calling for complex-systems thinking that
links processes across time scales, long-term experiments on
physical--chemical--biological interactions, attention to connectivity and
pattern across spatial scales, and multidisciplinary consortia that share
data and terminology under known boundary conditions. The binding
constraint that emerges from these appeals is not a single missing
experiment but a coupled deficit: coordinated, long-term, spatially
replicated, and harmonised data on the one hand, and better process
representation in models on the other, must advance together, because
neither longer records nor more elaborate models resolve the field's open
questions alone.

This paper is the research-frontier component of a larger review of tile
drainage, and it is deliberately framed around uncertainty. Where the
companion component consolidates settled knowledge, our purpose here is to
gather, structure, and sharpen the open questions---the priorities, the
hypotheses that remain to be tested, and the methods, instrumentation, and
datasets still needed to test them. We therefore do not attempt to resolve
any of the gaps we identify; we set them up, tracing how each is posed in
the primary literature and how the strands interlock. The aim is a map of
where the field should go next, in which the confidence of a claim about
what is \emph{not} yet understood is treated as a finding in its own right.

The synthesis proceeds from enabling frontiers to domain frontiers and
then to their consequences and integration. We begin with the cross-cutting
foundations on which every later domain depends: the observational
foundation of longer, denser, and harmonised monitoring; the process-based
drainage and water-quality models whose credibility is limited less by
their formulation than by the narrowness of the conditions under which they
have been tested; and the linked problems of mapping the largely unmapped
physical drainage network and scaling instrumented-field knowledge to the
watersheds where water-quality targets are actually set. On this footing we
turn to the domain frontiers---managing nitrogen, phosphorus, and
cropping-system losses; active water-table management and the re-evaluation
of drainage-system design; edge-of-field treatment technologies such as
denitrifying bioreactors, wetlands, and phosphorus-removal structures; and
salinity reclamation, drainage-water reuse, and disposal in arid and
irrigated systems. Because removing a constituent from one pathway can
shift it to another and altering the field water regime can set slow
biogeochemical and pedological change in motion, we then consolidate the
unintended consequences and multi-objective tradeoffs---greenhouse-gas
coupling, pollution swapping, emerging contaminants, trace-element
legacies, and drainage-induced soil evolution---that a nitrate- and
yield-centred accounting has historically overlooked. Finally, we draw the
integrative strands together, because translating frontier science into
outcomes depends on how practices are bundled and spatially targeted, how
they perform under a changing climate, and whether they are economically
and institutionally adoptable.

A single thread runs through all of these: no individual practice or model
is a silver bullet, and the decisive advances will come from coupling
mechanistic understanding with spatial targeting, economics, policy, and
shared decision-support infrastructure. Underlying that ambition is a data
imperative---better, longer, and more spatially representative
datasets---that recurs as the foundational requirement for validating
models, resolving tradeoffs, and delivering site-specific, climate-resilient
guidance \parencite{Ma2025,SKAGGS1994,Deng2011}. The sections that follow
elaborate that agenda frontier by frontier.